\begin{textsample}{POS Dim 3 – ap – Score 15.00 – ap/ap\_inf\_8.txt}  \label{ex:f3_pos_143}
3.4.
Direito de Aprendizagem e Desenvolvimento na Educação \textbf{Infantil} e Campos de Experiências A BNCC orienta que na etapa da Educação \textbf{Infantil} sejam assegurados seis Direitos de aprendizagens, considerando : as formas pelas quais \textbf{bebês}, \textbf{crianças} bem \textbf{pequenas} e \textbf{crianças} \textbf{pequenas} aprendem e constroem significações sobre si, os outros e o mundo social e natural e as exigências fundamentais da vida contemporânea ( BRASIL, 2017 ).
Os seis Direitos de Aprendizagens e Desenvolvimento são : I - Conviver com outras \textbf{crianças} e \textbf{adultos}, em \textbf{pequenos} e grandes grupos, utilizando diferentes linguagens, ampliando o conhecimento de si e do outro, o respeito em relação à cultura e às diferenças entre as pessoas.
II - \textbf{Brincar} cotidianamente de diversas formas, em diferentes espaços e tempos, com diferentes parceiros ( \textbf{crianças} e \textbf{adultos} ), ampliando e diversificando seu acesso a produções culturais, seus conhecimentos, sua imaginação, sua criatividade, suas experiências emocionais, corporais, \textbf{sensoriais}, expressivas, cognitivas, sociais e relacionais.
III - Participar ativamente, com \textbf{adultos} e outras \textbf{crianças}, tanto do planejamento da gestão da escola e das atividades propostas pelo educador quanto da realização das atividades da vida cotidiana, tais como a escolha das \textbf{brincadeiras}, dos materiais e dos ambientes, desenvolvendo diferentes linguagens e elaborando conhecimentos, decidindo e se posicionando.
IV - Explorar movimentos, \textbf{gestos}, \textbf{sons}, formas, \textbf{texturas}, cores, palavras, emoções, transformações, relacionamentos, histórias, objetos, elementos da natureza, na escola e fora dela, ampliando seus saberes sobre a cultura, em suas diversas modalidades : as artes, a escrita, a ciência e a tecnologia.
V - Expressar, como sujeito dialógico, criativo e sensível, suas necessidades, emoções, sentimentos, dúvidas, hipóteses, descobertas, opiniões, questionamentos, por meio de diferentes linguagens.
VI - Conhecer -se e construir sua identidade pessoal, social e cultural, constituindo uma imagem positiva de si e de seus grupos de pertencimento, nas diversas experiências de \textbf{cuidados}, interações, \textbf{brincadeiras} e linguagens vivenciadas na instituição escolar e em seu contexto familiar e comunitário. ( BRASIL, 2017 ).
Assim, os Campos e Experiências são pensados para assegurar esses direitos, sua organização curricular possui caráter interdisciplinar, conservando relações com as áreas de conhecimento visando à promoção do desenvolvimento cognitivo da \textbf{criança}.
Os cinco Campos de Experiências são : - O Eu, o Outro e o Nós ; - Corpo, \textbf{Gestos} e Movimentos ; - Traços, \textbf{Sons}, Cores e Formas ; - \textbf{Escuta}, Fala, Pensamento e Imaginação e - Espaços, tempos, quantidades, relações e transformações. ( BRASIL, 2017 ).
Na sequência são apresentados os Direitos de Aprendizagem e Desenvolvimento propostos para a Educação \textbf{Infantil}, sendo situados em cada campo de experiências ( BRASIL, 2017 ) e ajustados ao foco básico dos mesmos, dentro da divisão apresentada pela BNCC.

% matched lemmas: adulto, bebê, brincadeira, brincar, criança, cuidado, escuta, gesto, infantil, pequeno, sensorial, som, textura
\end{textsample}
